\section*{Блок 3. Гауссовские векторы и их свойства}
\addcontentsline{toc}{section}{Блок 3. Гауссовские векторы и их свойства}

\subsection*{1. Семантический анализ и математическая генерализация}
\addcontentsline{toc}{subsection}{1. Семантический анализ и математическая генерализация}

Многомерное нормальное распределение представляет собой фундаментальный объект классической теории вероятностей. Ключевое аналитическое свойство гауссовских векторов заключается в том, что для них \textbf{некоррелированность эквивалентна независимости}, а класс нормальных распределений замкнут относительно линейных преобразований.

Анализ задач данного блока выявляет два основных вычислительных паттерна:
\begin{enumerate}
    \item \textbf{Ортогональное проецирование:} Сведение задачи о вычислении условных математических ожиданий от зависимых компонент к работе с независимыми случайными величинами (Теорема о нормальной корреляции).
    \item \textbf{Вычисление моментов и нелинейных функционалов:} Использование аппарата производящих и характеристических функций, а также комбинаторных тождеств для многочленов от гауссовских величин.
\end{enumerate}

Генерализированная постановка опорной задачи формулируется следующим образом: дан центрированный гауссовский вектор $(X, Y, Z) \sim \N(0, \Sigma)$. Требуется вычислить нелинейное условное математическое ожидание $\E(f(Y,Z) | X)$. Прямое интегрирование с использованием многомерных условных плотностей вычислительно неэффективно. Оптимальный подход базируется на алгебраическом разложении пространства случайных величин.

\subsection*{2. Фундаментальный инструментарий: Условные распределения}
\addcontentsline{toc}{subsection}{2. Фундаментальный инструментарий: Условные распределения}

Для произвольного гауссовского вектора условное математическое ожидание одной подсистемы относительно другой всегда представляет собой линейную функцию, а условная ковариационная матрица детерминирована (не зависит от условия).

\begin{MethodBox}[Теорема о нормальной корреляции (Проекция)]
Пусть $(X, Y)$ — компоненты центрированного гауссовского вектора. Случайную величину $Y$ всегда можно представить в виде ортогональной суммы:
\[
    Y = aX + Y^\perp, \quad \text{где } Y^\perp \indep X
\]
Коэффициент регрессии $a$ определяется из условия ортогональности (некоррелированности) остатка $Y^\perp$ и базиса $X$:
\[
    \Cov(Y^\perp, X) = 0 \implies \Cov(Y - aX, X) = 0 \implies a = \frac{\Cov(X,Y)}{\Var(X)}.
\]
Дисперсия ортогонального остатка (она же условная дисперсия $\Var(Y|X)$) вычисляется как дополнение Шура:
\[
    \Var(Y^\perp) = \Var(Y) - \frac{\Cov^2(X,Y)}{\Var(X)}.
\]
\end{MethodBox}

\subsection*{3. Алгоритм метода ортогональных проекций}
\addcontentsline{toc}{subsection}{3. Алгоритм метода ортогональных проекций}

Ниже приведен строгий алгоритм редукции сложного условного математического ожидания к безусловному, применимый к любым функциям $f(Y,Z)$.

\begin{MethodBox}
\begin{enumerate}[label=\textbf{\arabic*.}]
    \item \textbf{Линейная параметризация.} Осуществляется переход к ортогональным остаткам:
    \[
        Y = aX + \tilde{Y}, \quad Z = bX + \tilde{Z}, \quad \text{где } a = \frac{\Cov(Y,X)}{\Var(X)}, \quad b = \frac{\Cov(Z,X)}{\Var(X)}.
    \]
    \item \textbf{Анализ ковариационной структуры остатков.} По построению $\tilde{Y} \indep X$ и $\tilde{Z} \indep X$. Вычисляется их взаимная ковариация:
    \[
        \Cov(\tilde{Y}, \tilde{Z}) = \Cov(Y - aX, Z - bX) = \Cov(Y,Z) - a b \Var(X).
    \]
    \textit{Следствие:} В подавляющем большинстве учебных задач матрица $\Sigma$ конструируется таким образом, что $\Cov(\tilde{Y}, \tilde{Z}) = 0$. В силу гауссовости это гарантирует независимость $\tilde{Y} \indep \tilde{Z}$.
    
    \item \textbf{Редукция условного ожидания.} С учетом независимости остатков от условия $X$, переменная $X$ внутри оператора $\E(\cdot|X)$ фиксируется как константа $x$, а условное ожидание заменяется безусловным:
    \[
        \E(f(Y,Z)|X=x) = \E(f(ax + \tilde{Y}, bx + \tilde{Z})).
    \]
    Дальнейшее вычисление сводится к интегралам по независимым стандартным распределениям.
\end{enumerate}
\end{MethodBox}

\subsection*{4. Продвинутые аналитические методы}
\addcontentsline{toc}{subsection}{4. Продвинутые аналитические методы}

Для обхода многомерного интегрирования при вычислении моментов высоких порядков и экспоненциальных функционалов применяется следующий аналитический аппарат.

\begin{HackBox}[Теорема Иссерлиса (Формула Вика)]
Для любого центрированного гауссовского вектора математическое ожидание произведения нечетного числа его компонент тождественно равно нулю. Для произведения четного числа (например, четырех) компонент справедлива формула суммирования по всем паросочетаниям:
\[
    \E(X_1 X_2 X_3 X_4) = \E(X_1X_2)\E(X_3X_4) + \E(X_1X_3)\E(X_2X_4) + \E(X_1X_4)\E(X_2X_3)
\]
\textit{Пример:} $\E(X^2 Y Z) = \E(X \cdot X \cdot Y \cdot Z) = \Var(X)\Cov(Y,Z) + 2\Cov(X,Y)\Cov(X,Z)$. Результат получается прямой подстановкой элементов матрицы $\Sigma$.
\end{HackBox}

\begin{HackBox}[Многомерная производящая функция моментов]
Вычисление математического ожидания экспоненты от линейной комбинации $\vec{t}^T \vec{\xi}$, где $\vec{\xi} \sim \N(\vec{\mu}, \Sigma)$, не требует интегрирования. Применяется тождество:
\[
    \E\left[ \exp(\vec{t}^T \vec{\xi}) \right] = \exp\left( \vec{t}^T \vec{\mu} + \frac{1}{2} \vec{t}^T \Sigma \vec{t} \right)
\]
\textit{Пример:} Для $\E(e^{X+Y+Z})$ вектор коэффициентов $\vec{t} = (1, 1, 1)^T$. Результат равен экспоненте от половины суммы всех элементов ковариационной матрицы $\Sigma$.
\end{HackBox}

\begin{HackBox}[Формула Шеппарда (Вероятности ортантов)]
Для центрированного двумерного гауссовского вектора $(X, Y)$ с коэффициентом корреляции $\rho = \frac{\Cov(X,Y)}{\sqrt{\Var(X)\Var(Y)}}$ вероятность попадания в первую четверть вычисляется аналитически:
\[
    \P(X > 0, Y > 0) = \frac{1}{4} + \frac{1}{2\pi}\arcsin(\rho)
\]
Из соображений симметрии: $\P(XY < 0) = \P(X>0, Y<0) + \P(X<0, Y>0) = \frac{1}{\pi}\arccos(\rho)$.
\end{HackBox}

\begin{HackBox}[Характеристическая функция для тригонометрических функционалов]
Если требуется вычислить математическое ожидание косинуса или синуса от гауссовской величины $W \sim \N(\mu, \sigma^2)$, используется характеристическая функция $\varphi_W(t) = \E(e^{itW}) = e^{it\mu - \frac{1}{2}\sigma^2 t^2}$.
В силу формулы Эйлера $\cos W = \Re(e^{iW})$, следовательно (при $t=1$):
\[
    \E(\cos W) = \exp\left(-\frac{\sigma^2}{2}\right) \cos(\mu)
\]
\end{HackBox}

\subsection*{5. Методы верификации и типичные ошибки}
\addcontentsline{toc}{subsection}{5. Методы верификации и типичные ошибки}

Для предотвращения фатальных ошибок на контрольной работе необходимо применять следующие методы самопроверки:

\begin{enumerate}
    \item \textbf{Проверка невырожденности.} Условная дисперсия $\Var(Y|X) = \Var(Y) - \frac{\Cov^2(X,Y)}{\Var(X)}$ обязана быть неотрицательной. Более того, она всегда меньше или равна безусловной дисперсии $\Var(Y)$. Получение условной дисперсии, превышающей исходную — стопроцентный маркер арифметической ошибки в знаке.
    
    \item \textbf{Анализ размерности (для формулы Вика).} Если исходное выражение имеет порядок $k$ (например, $X^2 Y Z$ — порядок 4), то итоговое выражение должно состоять из слагаемых, представляющих собой произведения $k/2$ ковариаций. Появление одиночных ковариаций в сумме с их произведениями недопустимо.
    
    \item \textbf{Симметрия ортогональных остатков.} После разложения $Y = aX + \tilde{Y}$ обязательно следует проверить математическое ожидание остатка. Поскольку $X$ и $Y$ центрированы, $\E(\tilde{Y})$ должно быть строго равно нулю.
\end{enumerate}

\begin{WarningBox}[Смещение вектора средних]
Представленный инструментарий (Формула Вика, ортогональные проекции через $\Cov$) применим \textbf{исключительно} к центрированным случайным векторам ($\vec{\mu} = \vec{0}$). 

Если вектор $(X, Y, Z)$ имеет ненулевое математическое ожидание $\vec{\mu} = (\mu_X, \mu_Y, \mu_Z)^T$, необходима предварительная замена переменных: $X_0 = X - \mu_X$, $Y_0 = Y - \mu_Y$. Все структурные разложения проводятся для $X_0, Y_0$, после чего осуществляется обратный переход. Игнорирование вектора средних полностью искажает ковариационную структуру проекций.
\end{WarningBox}
